\section{Experiments}
\label{sec:experiments}

\subsection{Datasets and protocol}
% TODO: fill in exact clip/frame counts once scripts/prepare_dataset.py
% has been run; see README.md "Dataset preparation".
We evaluate on three standard unsupervised VAD benchmarks: UCSD
Ped2, CUHK Avenue, and ShanghaiTech. All three provide normal-only
training video and mixed normal/anomalous test video with frame-level
(and, for Ped2, pixel-level) ground truth. We report frame-level ROC-AUC
and equal error rate (EER) using the standard per-clip min-max score
normalization protocol \cite{liu2018future}.

\subsection{Implementation details}
% TODO: report actual backbone/hyperparameters used for final numbers —
% these should match whatever is in configs/*.yaml at the time results
% were produced; keep this section in sync with the config, not the other
% way around.
Frame embeddings come from a frozen backbone (ResNet-18, $D{=}512$, or
DINOv2 ViT-S/14, $D{=}384$ — see \texttt{configs/*.yaml}). The causal SSM
stack has 2 layers with state size $N{=}128$ per layer
(Section~\ref{sec:method-ssm}), trained with AdamW, learning rate
$3\times10^{-4}$, on causal windows of 16 frames. All training and
evaluation was run on an Apple M3 Pro (MPS backend); no CUDA-specific
code path exists in this implementation.

\subsection{Main results}
Table~\ref{tab:main} reports frame-AUC, EER, and per-frame latency/FPS
measured directly on-device (Section~\ref{sec:method}, streaming
inference — \texttt{src/evaluate.py} calls the identical \texttt{step()}
routine used at deployment, with \texttt{torch.mps.synchronize()} before
each timing boundary). Rows attributed to prior work are as reported in
their respective papers, not reproduced under this protocol unless noted;
see \S\ref{sec:related} for citations.
% TODO: reproduce at least one baseline (VADMamba has public code) under
% this codebase's exact evaluation protocol before submission — a
% same-protocol comparison is much stronger evidence than citing numbers
% computed under a different eval setup.

% Auto-generated by scripts/make_latex_table.py — do not edit by hand.
\begin{table*}[t]
\centering
\caption{Main results: standard accuracy metrics alongside on-device streaming latency/throughput.}
\label{tab:main}
\begin{tabular}{llccccc}
\toprule
Dataset & Method & Frame-AUC$\uparrow$ & EER$\downarrow$ & Latency (ms)$\downarrow$ & FPS$\uparrow$ & Hardware \\
\midrule
UCSD Ped2 & VADMamba (ICME'25) & -- & -- & -- & -- & -- \\ % source: arXiv:2503.21169
UCSD Ped2 & Ours (streaming SSM) & 67.9 & 36.8 & 0.74 & 1343.62 & mps \\
\midrule
CUHK Avenue & VADMamba (ICME'25) & -- & -- & -- & -- & -- \\ % source: arXiv:2503.21169
CUHK Avenue & Ours (streaming SSM) & 70.2 & 34.9 & 0.77 & 1290.74 & mps \\
\bottomrule
\end{tabular}
\end{table*}


\subsection{Theory validation}
Table~\ref{tab:theory} compares the settling-delay bound
(Eq.~\ref{eq:delay-bound}, Proposition~\ref{prop:settling}) computed from
each trained layer's mean learned base decay against the empirically
measured detection delay (frames between an anomalous segment's labeled
onset and the first score-threshold crossing;
\texttt{src.metrics.detection\_delay}).

% Auto-generated by scripts/make_theory_table.py — do not edit by hand.
\begin{table*}[t]
\centering
\caption{Predicted settling delay (Eq.~\ref{eq:delay-bound}), from each dataset's trained mean decay, vs.\ empirically measured detection delay.}
\label{tab:theory}
\resizebox{\linewidth}{!}{%
\small
\begin{tabular}{llccc}
\toprule
Dataset & Layer & Mean decay $a$ & Theory delay (frames) & Empirical delay (frames) \\
\midrule
UCSD Ped2 & 0 & 0.9493 & 57.5 & 1.6 ($n{=}$12) \\
 & 1 & 0.9501 & 58.6 &  \\
\midrule
CUHK Avenue & 0 & 0.9493 & 57.5 & 18.4 ($n{=}$21) \\
 & 1 & 0.9501 & 58.6 &  \\
\bottomrule
\end{tabular}
}
\end{table*}


\subsection{Ablations}
Tables~\ref{tab:ablation-ucsd-ped2} and \ref{tab:ablation-cuhk-avenue}
report decay-rate, state-size, and event-boundary-gate ablations on Ped2
and Avenue respectively, all else (backbone, training budget, epochs)
held fixed at the Section~\ref{sec:experiments} implementation details.

% Auto-generated by scripts/make_ablation_table.py — do not edit by hand.
\begin{table*}[t]
\centering
\caption{Ablations on UCSD Ped2: architectural variants, all else held fixed. Latency/FPS measured the same way as Table~\ref{tab:main}.}
\label{tab:ablation-ucsd-ped2}
\begin{tabular}{lcccc}
\toprule
Variant & Frame-AUC$\uparrow$ & EER$\downarrow$ & Latency (ms)$\downarrow$ & FPS$\uparrow$ \\
\midrule
Baseline (gate on, N=128, decay in [.9,.999]) & 67.9 & 36.8 & 0.74 & 1344 \\
Gate off & 76.6 & 24.6 & 0.49 & 2047 \\
Decay fast, decay in [.5,.9] & 67.0 & 38.2 & 0.75 & 1330 \\
Decay slow, decay in [.98,.9999] & 67.8 & 37.0 & 0.76 & 1311 \\
State size N=32 & 74.7 & 33.1 & 0.77 & 1302 \\
State size N=256 & 66.5 & 38.7 & 0.77 & 1298 \\
\bottomrule
\end{tabular}
\end{table*}

% Auto-generated by scripts/make_ablation_table.py — do not edit by hand.
\begin{table*}[t]
\centering
\caption{Ablations on CUHK Avenue: architectural variants, all else held fixed. Latency/FPS measured the same way as Table~\ref{tab:main}.}
\label{tab:ablation-cuhk-avenue}
\begin{tabular}{lcccc}
\toprule
Variant & Frame-AUC$\uparrow$ & EER$\downarrow$ & Latency (ms)$\downarrow$ & FPS$\uparrow$ \\
\midrule
Baseline (gate on, N=128, decay in [.9,.999]) & 70.2 & 34.9 & 0.77 & 1291 \\
Gate off & 60.0 & 43.8 & 0.51 & 1953 \\
Decay fast, decay in [.5,.9] & 70.6 & 34.4 & 0.75 & 1327 \\
Decay slow, decay in [.98,.9999] & 70.3 & 34.8 & 0.77 & 1303 \\
State size N=32 & 70.3 & 35.0 & 0.76 & 1319 \\
State size N=256 & 71.1 & 34.5 & 0.77 & 1302 \\
\bottomrule
\end{tabular}
\end{table*}


% TODO: this cross-dataset flip is the most interesting empirical result
% in the paper so far and deserves a dedicated discussion, not just this
% paragraph — consider a figure plotting AUC-vs-capacity per dataset once
% ShanghaiTech (a third, larger data point) is in.
The two datasets tell different stories, and the difference is
informative. On Ped2 (16 training clips, 120-180 frames each),
\emph{disabling} the gate improves frame-AUC over the gated baseline
(76.6\% vs.\ 67.9\%) and the smaller state ($N{=}32$) beats both the
baseline ($N{=}128$) and the larger variant ($N{=}256$) — i.e.\ every
higher-capacity variant underperforms a simpler one. On Avenue (also 16
training clips, but 36-1271 frames each — substantially more total
training frames despite the same clip count), the pattern
\emph{reverses}: disabling the gate \emph{hurts} sharply (60.0\% vs.\
70.2\%), and the larger state ($N{=}256$) slightly beats the baseline
rather than underperforming it. This is consistent with the gate's extra
parameters (Eq.~\ref{eq:gate}) overfitting on Ped2's very small training
set but finding real signal once given Avenue's larger one — evidence
for a capacity/data-size interaction rather than the gating mechanism
being unhelpful in principle (Section~\ref{sec:theory-gating}), but we
only have two data points and both happen to have 16 training clips (only
total frame count differs) — a third, larger dataset (ShanghaiTech) is
needed before treating this as more than a hypothesis. The decay-rate
sweep (fast vs.\ slow base decay, gate on) shows comparatively little
effect on either dataset, suggesting decay range is not the dominant
factor. We do not yet have a strict-causal-vs.-windowed ablation (would
require a lookahead-buffer variant of Section~\ref{sec:method-ssm}, not
yet implemented).

\subsection{Edge deployment}
% TODO: CoreML/ANE export + benchmark not yet automated in this
% codebase (see README.md "Reproducing the edge-latency numbers");
% MPS latency numbers in Table~\ref{tab:main} are already real
% on-device measurements, not simulated.
